\documentclass[11pt]{article}
\usepackage[T1]{fontenc}
\usepackage{amsmath,amssymb,amsthm}
\usepackage[margin=1in]{geometry}
\usepackage{booktabs}
\usepackage[numbers]{natbib}
\usepackage[hidelinks]{hyperref}
\hypersetup{
  pdftitle={Exact Facial Distance and the Minimizing Face Orbit of the 3 x 3 Birkhoff Polytope},
  pdfauthor={Xiaojian Zeng},
  pdfsubject={Exact facial distance and face-orbit classification},
  pdfkeywords={Birkhoff polytope, facial distance, face orbit, exact certificate}
}

\newtheorem{theorem}{Theorem}
\newtheorem{lemma}[theorem]{Lemma}
\DeclareMathOperator{\dist}{dist}
\DeclareMathOperator{\aff}{aff}
\DeclareMathOperator{\conv}{conv}
\DeclareMathOperator{\vertices}{vert}

\title{Exact Facial Distance and the Minimizing Face Orbit of the
$3\times3$ Birkhoff Polytope}
\author{Xiaojian Zeng\\
PBC School of Finance, Tsinghua University, 100190\\
\texttt{xiaojianzeng@foxmail.com}}
\date{}

\begin{document}
\maketitle

\begin{abstract}
For a nonempty proper face $G$ of the $3\times3$ Birkhoff polytope
$\mathcal B_3$, let $\mathcal B_3[G]$ be the convex hull of the permutation
vertices outside $G$.  With the Frobenius norm, we prove
\[
 \min_G \dist\bigl(G,\mathcal B_3[G]\bigr)=\frac1{\sqrt2}.
\]
All minimizing faces are edges, and they form one orbit under independent
row and column permutations and transposition: the six edges whose two
permutation vertices have relative permutation a $3$-cycle.  The proof
enumerates all $48$ nonempty proper faces into five symmetry orbits and gives
exact rational projection and simplex optimality certificates for one
representative of each orbit.  The statement is specific to $\mathcal B_3$;
no general-dimensional or full-automorphism-group claim is made.
\end{abstract}

\section{Statement and context}

Write $S_3$ in one-line notation, and let $P_\sigma$ be the permutation
matrix with $(P_\sigma)_{i,\sigma(i)}=1$.  The Birkhoff polytope is
\[
 \mathcal B_3=\left\{X\in\mathbb R_{\geq0}^{3\times3}:
 X\mathbf1=\mathbf1,\ X^{\mathsf T}\mathbf1=\mathbf1\right\}
 =\conv\{P_\sigma:\sigma\in S_3\}.
\]
Its facial combinatorics belongs to the broader study of faces of Birkhoff
polytopes; see, for example, \citet{Paffenholz2015Faces}.  For a nonempty
proper face $G$, put
\[
 \mathcal B_3[G]=\conv\bigl(\vertices(\mathcal B_3)\setminus G\bigr),
 \qquad
 d(G)=\dist\bigl(G,\mathcal B_3[G]\bigr),
\]
where distance is measured in the Frobenius norm, and define
\[
 \Phi(\mathcal B_3)=\min_{\substack{\varnothing\ne G\subsetneq\mathcal B_3\\
 G\text{ a face}}}d(G).
\]
This is the complement-hull formulation used by \citet{PenaWirth2026}.
Their Example~6.4 gives the general lower bound
$\Phi(\mathcal B_n)\geq 1/\sqrt{n^2-1}$, which becomes $1/\sqrt8$ for
$n=3$, but does not compute the exact value in this case.

Let $H$ be the group generated by independent row and column permutations
and transposition.  Its elements act as
$X\mapsto RXP^{\mathsf T}$ or $X\mapsto RX^{\mathsf T}P^{\mathsf T}$,
where $R$ and $P$ are permutation matrices.  These give $72$ distinct
actions on the six vertices.  We classify equality only under this specified
group; we do not assert that $H$ is the full affine-isometry group.

\begin{theorem}\label{thm:main}
One has
\[
 \Phi(\mathcal B_3)=\frac1{\sqrt2}.
\]
Exactly six faces attain this value.  They form a single $H$-orbit,
represented by
\[
 G=\conv\{P_{123},P_{231}\};
\]
equivalently, the two endpoints of every minimizing edge have relative
permutation a $3$-cycle.
\end{theorem}

The finite-certificate design is related in spirit to the exact
symmetry-reduced enumeration for finite problems in discrete geometry
documented in \cite{ZengEtAl2026FiniteCertificates}.  Here the combinatorial
enumeration and the metric certificates are kept separate.  This comparison
concerns verification architecture only, not a transfer of a mathematical
result.  No literature-priority claim is made for Theorem~\ref{thm:main}.

\section{Complete face and orbit reduction}

For $S\subseteq S_3$, define the union of its matrix supports by
\[
 E(S)=\bigcup_{\sigma\in S}\{(i,\sigma(i)):i=1,2,3\}
\]
and its support closure by
\[
 \operatorname{cl}(S)=
 \{\tau\in S_3:\{(i,\tau(i)):i=1,2,3\}\subseteq E(S)\}.
\]

\begin{lemma}\label{lem:closure}
A nonempty proper subset $S\subset S_3$ is the vertex set of a face of
$\mathcal B_3$ if and only if $\operatorname{cl}(S)=S$.
\end{lemma}

\begin{proof}
Inside its affine hull, $\mathcal B_3$ is defined by the nine coordinate
inequalities $x_{ij}\geq0$.  The smallest face containing the vertices
$\{P_\sigma:\sigma\in S\}$ is obtained by setting $x_{ij}=0$ at every
coordinate absent from $E(S)$.  Its permutation vertices are therefore
exactly those indexed by $\operatorname{cl}(S)$.  Conversely, the common
zero coordinates of a face expose precisely this support-closed vertex set.
\end{proof}

Applying Lemma~\ref{lem:closure} to the $2^6$ subsets is completely
explicit.  All $6$ singletons and all $15$ pairs are closed.  Among the
$20$ triples, only the two parity classes
\[
 \{123,231,312\},\qquad \{132,213,321\}
\]
are not closed; each uses all nine coordinates and hence closes to $S_3$.
The closed four-subsets are exactly the nine sets
\[
 \{\sigma\in S_3:\sigma(i)\ne j\},\qquad 1\leq i,j\leq3,
\]
and no five-subset is closed.  Direct affine-rank calculation gives the
dimensions in Table~\ref{tab:orbits}.  Relative cycle type separates the two
edge orbits; row and column permutations and inversion by transposition are
transitive on each of the other displayed classes.  Thus the list has
$6+6+9+18+9=48$ faces and omits none.

\begin{table}[ht]
\centering
\caption{All nonempty proper faces modulo $H$.}
\label{tab:orbits}
\begin{tabular}{@{}ccc@{}}
\toprule
dimension & orbit size & representative vertex labels \\
\midrule
$0$ & $6$  & $\{123\}$ \\
$1$ & $6$  & $\{123,231\}$ \\
$1$ & $9$  & $\{123,132\}$ \\
$2$ & $18$ & $\{123,132,213\}$ \\
$3$ & $9$  & $\{123,132,213,231\}$ \\
\bottomrule
\end{tabular}
\end{table}

Every element of $H$ is a Frobenius isometry and permutes the six vertices.
It consequently preserves both $G\mapsto\mathcal B_3[G]$ and $d(G)$, so one
distance calculation per row of Table~\ref{tab:orbits} suffices.

\section{Exact distance certificates}

Let $G=\conv\{A_1,\ldots,A_r\}$ be one of the five representatives, and
list its complementary vertices as $C(G)=(B_1,\ldots,B_s)$.  We first
compute the relaxed distance
\[
 \delta(G)^2=
 \min_{\substack{\mathbf1^{\mathsf T}\alpha=1\\
                   \beta\geq0,\ \mathbf1^{\mathsf T}\beta=1}}
 \left\|\sum_{i=1}^r\alpha_iA_i-
             \sum_{j=1}^s\beta_jB_j\right\|_F^2
 =\dist\bigl(\aff G,\mathcal B_3[G]\bigr)^2.
\]
The affine coefficients $\alpha$ need not be nonnegative.

Set $x=\sum_i\alpha_iA_i$, $y=\sum_j\beta_jB_j$, and $e=x-y$.  A pair
$(\alpha,\beta)$ certifies the minimum if
\begin{align}
 \langle e,A_i-A_1\rangle_F&=0 &&(i=2,\ldots,r),\label{eq:orth}\\
 g_j:=-2\langle B_j,e\rangle_F&=\mu
     &&\text{when }\beta_j>0,\label{eq:active}\\
 g_j&\geq\mu &&\text{when }\beta_j=0.\label{eq:inactive}
\end{align}
Indeed, \eqref{eq:orth} makes $x$ the orthogonal projection of $y$ onto
$\aff G$.  After this projection, $g_j$ is the derivative with respect to
$\beta_j$ of the resulting convex quadratic.  Conditions
\eqref{eq:active}--\eqref{eq:inactive} are the simplex optimality
conditions, so they are sufficient for a global minimum.

Denote the representatives in Table~\ref{tab:orbits} from top to bottom by
$G_0,\ldots,G_4$.  The complementary vertices, in the coordinate order used
below, are
\begin{align*}
 C(G_0)&=(132,213,231,312,321),&
 C(G_1)&=(132,213,312,321),\\
 C(G_2)&=(213,231,312,321),&
 C(G_3)&=(231,312,321),\\
 C(G_4)&=(312,321).&&
\end{align*}
Table~\ref{tab:certificates} gives rational certificates.  In each row the
entries of $\alpha$ sum to one, $\beta$ is a probability vector, and direct
substitution into the six permutation matrices makes every inner product in
\eqref{eq:orth} zero.  The displayed vector $g$ verifies
\eqref{eq:active}--\eqref{eq:inactive}; its common active value is its
smallest entry.  The final column is the exact value of
$\langle e,e\rangle_F$.

\begin{table}[ht]
\centering
\caption{Exact affine-distance and simplex optimality certificates.}
\label{tab:certificates}
\setlength{\tabcolsep}{3.5pt}
\small
\begin{tabular}{@{}ccccc@{}}
\toprule
face & $\alpha$ & $\beta$ & $g$ & $\delta(G)^2$ \\
\midrule
$G_0$ & $(1)$
& $(\frac13,\frac13,0,0,\frac13)$
& $(0,0,2,2,0)$ & $2$ \\
$G_1$ & $(\frac12,\frac12)$
& $(\frac13,\frac13,0,\frac13)$
& $(0,0,2,0)$ & $\frac12$ \\
$G_2$ & $(\frac12,\frac12)$
& $(\frac14,\frac14,\frac14,\frac14)$
& $(\frac32,\frac32,\frac32,\frac32)$ & $\frac94$ \\
$G_3$ & $(-\frac15,\frac35,\frac35)$
& $(\frac12,\frac12,0)$
& $(\frac35,\frac35,\frac{12}5)$ & $\frac9{10}$ \\
$G_4$ & $(\frac14,\frac14,\frac14,\frac14)$
& $(\frac12,\frac12)$
& $(3,3)$ & $\frac94$ \\
\bottomrule
\end{tabular}
\end{table}

We now pass from these relaxed distances to the required facial distances.
Because $G\subseteq\aff G$,
\[
 d(G)\geq\delta(G)
\]
for every face.  Table~\ref{tab:certificates} therefore gives
$d(G)\geq1/\sqrt2$ for all $48$ nonempty proper faces.  For $G_1$, the
affine vector $\alpha=(1/2,1/2)$ is already convex.  More explicitly,
\[
 x=\frac{P_{123}+P_{231}}2\in G_1,
 \qquad
 y=\frac{P_{132}+P_{213}+P_{321}}3
   =\frac13\mathbf1\mathbf1^{\mathsf T}\in\mathcal B_3[G_1].
\]
The matrix $x-y$ has six entries equal to $1/6$ and three entries equal to
$-1/3$, so
\[
 \|x-y\|_F^2=6\left(\frac16\right)^2+
                  3\left(\frac13\right)^2=\frac12.
\]
Thus $d(G_1)=1/\sqrt2$.  The other four rows have
$\delta(G)^2>1/2$, so no face in any other orbit can attain equality.  The
orbit of $G_1$ has size six, proving Theorem~\ref{thm:main}.

\section{Computational and AI assistance}

The face closure and the five tables were checked with deterministic Python
programs supplied as ancillary material.  One program enumerates the six
permutation vertices, all support closures, and all $72$ specified symmetry
actions using integer set operations.  A second program reconstructs every
entry of Table~\ref{tab:certificates} with exact rational arithmetic and
checks coefficient sums, projection orthogonality, simplex feasibility,
gradient inequalities, squared norms, and the unique minimum.  No
floating-point tolerance or randomized step is used.  The proof above
includes the complete finite classification and the rational certificates,
so the programs serve as reproducibility checks rather than numerical
evidence in place of a proof.

An AI coding assistant helped organize the argument, prepare the LaTeX, and
execute the deterministic checks.  A separate AI-assisted audit recomputed
the face classification and the rational optimality conditions.  These are
tool-assisted review passes, not claims of formal verification or independent
peer review; responsibility for the mathematical assertions remains with the
author.

\bibliographystyle{plainnat}
\bibliography{references}

\end{document}
